\section{Objetivos de investigación}
\label{sec:objectives}

El objetivo general de esta investigación es desarrollar un marco computacional multiescala matemáticamente riguroso que derive los parámetros de transporte ambiental macroscópico directamente a partir de la mecánica fluido--partícula microestructural con bajos números de Reynolds. Para alcanzar este propósito global, el proyecto cumplirá con los siguientes objetivos específicos:

\begin{enumerate}
    \item \textbf{Implementar un módulo de Dinámica Stokesiana de alta fidelidad:} Configurar y optimizar un motor de seguimiento hidrodinámico de muchos cuerpos de código abierto (por ejemplo, LAMMPS o librerías especializadas en SD) para simular las trayectorias de suspensiones de partículas rígidas ($Re \ll 1$) dentro de geometrías bentónicas complejas. Esto incluye incorporar explícitamente matrices de movilidad de muchos cuerpos de campo lejano, fuerzas de lubricación de corto rango y el tensor de velocidad de deformación del fluido ambiental ($\bm{{E}}^{\infty}$) bajo condiciones de cizallamiento localizado.
    \item \textbf{Caracterizar los mecanismos microestructurales de captura de partículas:} Cuantificar la eficiencia de captura discreta y la dinámica de acumulación de partículas coloidales que interactúan con interfaces de frontera complejas, mapeando específicamente las restricciones físicas de las arquitecturas de la cubierta coralina (por ejemplo, la intercepción en la capa límite y el estancamiento por vórtices de estela) y las gargantas de poro de sedimentos de alta permeabilidad características de las bahías del Parque Nacional Natural Tayrona \cite{goldman1967shear, monismith2007hydrodynamics, sen2010lightdriven}.
    \item \textbf{Desarrollar un protocolo determinista de escalamiento de micro a macro:} Establecer una interfaz de mapeo matemático para llevar los datos de configuración de partículas discretas hacia campos continuos. Este objetivo implica calcular el Tensor de Dispersión Hidrodinámica Efectiva ($\bm{D}_{\text{eff}}$) a través de la varianza asintótica a tiempos largos del desplazamiento de las partículas, y calcular la tasa de captura cinética continua ($R$) mediante los flujos de masa netos a través de las fronteras.
    \item \textbf{Validar frente a la literatura experimental empírica:} Ejecutar resolvedores numéricos para la ecuación continua de Advección--Difusión--Reacción (ADR) utilizando los parámetros derivados de forma determinista. Validar críticamente la precisión predictiva del marco computacional escalado contrastando directamente los campos de concentración numérica y las eficiencias de captura frente a conjuntos de datos empíricos e independientes disponibles en la literatura, incluyendo datos físicos de seguimiento en canales de flujo de cubiertas coralinas y mediciones locales de permeabilidad en arenas de arrecife \cite{mendrik2024microplastic, reidenbach2006boundary}.
\end{enumerate}
